\chapter{Technology}\label{ch:b1:technology}

Two bakeries make the same bread. One uses eight bakers and two ovens, the
other four bakers and four ovens, and each turns out a few hundred loaves a
day. Neither is wasteful: each gets as much bread as it can from what it
uses. They differ because they pay different prices for bakers and ovens.
Across whole economies the same pattern is large and regular: where labour
is dear, each hour of work is done with more capital. This chapter
describes what a firm can make from its inputs --- its technology --- in the
language of \cref{ch:b1:preferences-and-utility}: \omterm{def:b1:technology:isoquant}{isoquants} instead of
\omterm{def:b1:preferences-and-utility:indifference-curve}{indifference curves}, a rate of technical substitution instead of a rate of
substitution. The next chapter adds prices.

\begin{openingchart}{Capital per hour worked against labour compensation per
hour, 2019, 112 countries, logarithmic axes (2021 US dollars at purchasing
power parity), with the least-squares line of the logarithms.}{Penn World
Table 11.0: capital stock at current PPPs divided by total hours (persons
engaged times average annual hours); labour share times output-side GDP at
current PPPs divided by total hours}
\begin{tikzpicture}
\begin{loglogaxis}[omscatter, log ticks with fixed point, xtickten={-1,...,2}, ytickten={0,...,3}, xmin=0.03, xmax=100, ymin=2, ymax=1000,
  xlabel={Labour compensation per hour (dollars)}, ylabel={Capital per hour worked (dollars)},
  width=0.7\linewidth, height=0.5\linewidth]
  \addplot[only marks, mark=*, mark size=1.3pt, omDef] table[x=wph, y=kph, col sep=comma]
    {figdata/bachelor-1/09-technology/out/capital_wage_2019.csv};
  \addplot[no marks, sharp plot, thick, omThm] table[x=wph, y=fit, col sep=comma]
    {figdata/bachelor-1/09-technology/out/capital_wage_fit.csv};
  \node[font=\footnotesize, anchor=west] at (axis cs:0.06,400) {slope $0.94$};
\end{loglogaxis}
\end{tikzpicture}
\end{openingchart}

The 112 countries lie close to a straight line with a slope of $0.94$ (the
correlation of the logarithms is $0.91$): a country where an hour of work
costs ten times more uses about $10^{0.94} = 8.7$ times more capital per hour.
% figdata: capital_wage_stats
An hour of work in the United States was paid 45.86 dollars and used 256.58
dollars of capital; in India it was paid 3.97 dollars and used 27.98.
% figdata: capital_wage_2019
The chart compares whole economies, not one industry, and capital is
measured in money, not in ovens; but the pattern is the one the bakeries
show, and this chapter gives the tools to read it.

\begin{recall}
A \omterm{def:b1:technology:production-function}{production function} with one variable input, and the diminishing
\omterm{def:b1:technology:marginal-product}{marginal product} of that input, were met at school, as was the aggregate
\omterm{def:b1:technology:production-function}{production function} $Y = AK^\alpha L^{1-\alpha}$ of growth accounting. A partial
derivative holds the other variables fixed; the chain rule along a level
curve gives its slope (\cref{ch:b1:constrained-optimisation}). \omterm{def:b1:preferences-and-utility:indifference-curve}{Indifference
curves} and the \omterm{def:b1:preferences-and-utility:mrs}{marginal rate of substitution} are in
\cref{ch:b1:preferences-and-utility}.
\end{recall}

\section{Production functions and isoquants}\label{sec:b1:technology:isoquants}

A firm turns inputs into output. With two inputs, capital $K$ (machine
hours, ovens) and labour $L$ (hours of work), a \emph{production plan}
$(K, L, Y)$ lists the inputs used and the output made.

\begin{definition}[Production function]\label{def:b1:technology:production-function}
\index{production function}\index{technically efficient}
A production plan is \emph{technically efficient} if no other plan makes
more output with no more of any input. The \emph{production function}
$Y = f(K, L)$ gives the largest output that the inputs $(K, L)$ can make: it
lists the technically efficient plans.
\end{definition}

\begin{definition}[Isoquant]\label{def:b1:technology:isoquant}
\index{isoquant}
An \emph{isoquant} is the set of input bundles that make the same output
$Y_0$: $\{(K, L) : f(K, L) = Y_0\}$.
\end{definition}

\begin{definition}[Marginal product]\label{def:b1:technology:marginal-product}
\index{marginal product}
The \emph{marginal product} of an input is the partial derivative of the
\omterm{def:b1:technology:production-function}{production function} with respect to it, $f_L = \partial f/\partial L$ and
$f_K = \partial f/\partial K$: the extra output from one more unit of the input,
the other held fixed. It is \emph{diminishing} if it falls as the input
rises, $f_{LL} < 0$.
\end{definition}

Unlike utility, output is measured in loaves: a \omterm{def:b1:technology:production-function}{production function} is
cardinal, and doubling it means twice the bread. \omterm{def:b1:technology:isoquant}{Isoquants} play the role of
\omterm{def:b1:preferences-and-utility:indifference-curve}{indifference curves}, but each carries a number that matters.

\section{The marginal rate of technical substitution}\label{sec:b1:technology:mrts}

\begin{definition}[Marginal rate of technical substitution]\label{def:b1:technology:mrts}
\index{marginal rate of technical substitution}\index{MRTS}
The \emph{marginal rate of technical substitution} at $(K, L)$ is the amount
of capital that can replace one unit of labour with output unchanged:
minus the slope $\dd K/\dd L$ of the \omterm{def:b1:technology:isoquant}{isoquant}, with $L$ on the horizontal
axis. By the chain rule along $f(K, L) = Y_0$,
\[
\mathrm{MRTS} = -\frac{\dd K}{\dd L}\bigg|_{Y = Y_0} = \frac{f_L}{f_K}.
\]
\end{definition}

The proof is that of the \omterm{def:b1:preferences-and-utility:mrs}{marginal rate of substitution}:
$f_K\,\dd K + f_L\,\dd L = 0$ along the \omterm{def:b1:technology:isoquant}{isoquant}. For $Y = K^{1/2}L^{1/2}$,
$f_L/f_K = K/L$: at $(K, L) = (16, 4)$ one hour of work replaces four machine
hours; at $(4, 16)$ it replaces a quarter. \omterm{def:b1:technology:isoquant}{Isoquants} are \emph{convex} when
the \omterm{def:b1:technology:mrts}{MRTS} falls as labour replaces capital along them, as here. Convexity
says that mixed plans do at least as well as extreme ones; it holds for
most technologies observed, though some inputs (two kinds of fuel) are
close to \omterm{def:b1:preferences-and-utility:perfect}{perfect substitutes}.

\begin{omfigure}
\begin{tikzpicture}
\begin{axis}[omrecord, grid=none, xlabel={Labour $L$}, ylabel={Capital $K$}, xmin=0, xmax=20,
  ymin=0, ymax=20, width=0.5\linewidth, height=0.44\linewidth, samples=80]
  \addplot[demand, domain=0.8:20] {16/x};
  \addplot[demand, domain=1.8:20] {36/x};
  \addplot[demand, domain=3.2:20] {64/x};
  \addplot[guide, domain=0:8] {8 - x};
  \addplot[only marks, mark=*, black] coordinates {(4,4)};
  \node[font=\footnotesize, anchor=south west] at (axis cs:4,4) {$\mathrm{MRTS} = 1$};
  \node[font=\footnotesize, anchor=south west] at (axis cs:12,1.33) {$Y = 4$};
  \node[font=\footnotesize, anchor=south west] at (axis cs:12,3) {$Y = 6$};
  \node[font=\footnotesize, anchor=south west] at (axis cs:12,5.33) {$Y = 8$};
\end{axis}
\end{tikzpicture}

{\small \omterm{def:b1:technology:isoquant}{Isoquants} of $Y = K^{1/2}L^{1/2}$ for outputs 4, 6 and 8. At $(4, 4)$ on
the \omterm{def:b1:technology:isoquant}{isoquant} $Y = 4$ the slope is $-1$ (the tangent is drawn dotted): one hour of
labour replaces one machine hour.}
\end{omfigure}

Three families of technologies recur. They differ in how easily one input
replaces the other.

\begin{definition}[Cobb--Douglas production function]\label{def:b1:technology:cobb-douglas}
\index{Cobb--Douglas production function}
A \emph{Cobb--Douglas \omterm{def:b1:technology:production-function}{production function}} is $f(K, L) = AK^\alpha L^\beta$ with
$A, \alpha, \beta > 0$. Its \omterm{def:b1:technology:mrts}{MRTS} is $(\beta/\alpha)(K/L)$.
\end{definition}

\begin{definition}[Fixed-proportions production function]\label{def:b1:technology:leontief}
\index{fixed-proportions production function}\index{Leontief production function}
A \emph{fixed-proportions} (Leontief) \omterm{def:b1:technology:production-function}{production function} is
$f(K, L) = \min(K/a, L/b)$: each unit of output needs $a$ units of capital and
$b$ of labour, and an extra unit of either input alone adds nothing. Its
\omterm{def:b1:technology:isoquant}{isoquants} are right angles with corners on the ray $K/L = a/b$.
\end{definition}

\omterm{def:b1:preferences-and-utility:perfect}{Perfect substitutes}, $f = aK + bL$, have straight \omterm{def:b1:technology:isoquant}{isoquants} and a constant
\omterm{def:b1:technology:mrts}{MRTS} $b/a$.

\begin{omfigure}
\begin{tikzpicture}
\begin{groupplot}[group style={group size=3 by 1, horizontal sep=10mm}, omrecord, grid=none,
  xmin=0, xmax=10, ymin=0, ymax=10, width=0.345\linewidth, height=0.34\linewidth, samples=60,
  xlabel={$L$}, ylabel={$K$}]
\nextgroupplot[title={\small Cobb--Douglas}, width=0.345\linewidth]
  \addplot[demand, domain=0.45:10] {4/x};
  \addplot[demand, domain=0.9:10] {9/x};
\nextgroupplot[title={\small Fixed proportions}, width=0.345\linewidth]
  \addplot[demand] coordinates {(2,10) (2,2) (10,2)};
  \addplot[demand] coordinates {(4,10) (4,4) (10,4)};
  \addplot[guide, domain=0:10] {x};
\nextgroupplot[title={\small Perfect substitutes}, width=0.345\linewidth]
  \addplot[demand, domain=0:4] {4 - x};
  \addplot[demand, domain=0:7] {7 - x};
\end{groupplot}
\end{tikzpicture}

{\small Two \omterm{def:b1:technology:isoquant}{isoquants} of $K^{1/2}L^{1/2}$, of $\min(K, L)$ (corners on the ray
$K = L$, dotted) and of $K + L$.}
\end{omfigure}

\section{Returns to scale}\label{sec:b1:technology:scale}

What happens when every input is scaled up together?

\begin{definition}[Returns to scale]\label{def:b1:technology:returns-to-scale}
\index{returns to scale}\index{constant returns to scale}\index{increasing returns to scale}\index{decreasing returns to scale}
A technology has \emph{constant}, \emph{increasing} or \emph{decreasing
returns to scale} if, for every $t > 1$ and every $(K, L)$,
$f(tK, tL)$ is equal to, greater than or less than $tf(K, L)$. A function is
\emph{homogeneous of degree $k$} if $f(tK, tL) = t^kf(K, L)$ for all $t > 0$; it
has constant returns if $k = 1$, increasing if $k > 1$, decreasing if $k < 1$.
\end{definition}

The Cobb--Douglas function $AK^\alpha L^\beta$ is homogeneous of degree $\alpha + \beta$.
\omterm{def:b1:technology:returns-to-scale}{Returns to scale} need not be the same at every size: for
$f = K^{1/2}L^{1/2} - 2$, doubling $(4, 4)$ to $(8, 8)$ raises output from 2 to 6.
The \emph{local} \omterm{def:b1:technology:returns-to-scale}{returns to scale} at a point are measured by the elasticity
of scale, $(Kf_K + Lf_L)/f$, the percentage rise of output when all inputs
rise by one per cent; for a homogeneous function it equals $k$ everywhere,
by the next theorem.

\begin{theorem}[Euler's theorem]\label{thm:b1:technology:euler}
If $f$ is differentiable and homogeneous of degree $k$, then
$Kf_K(K, L) + Lf_L(K, L) = kf(K, L)$.
\end{theorem}

\begin{proof}
Differentiate $f(tK, tL) = t^kf(K, L)$ with respect to $t$:
$Kf_K(tK, tL) + Lf_L(tK, tL) = kt^{k-1}f(K, L)$. Set $t = 1$.
\end{proof}

With constant returns, paying each input its \omterm{def:b1:technology:marginal-product}{marginal product} uses up the
output exactly: $Kf_K + Lf_L = f$. \Cref{ch:b1:factor-markets} returns to
this.

\begin{omfigure}
\begin{tikzpicture}
\begin{groupplot}[group style={group size=2 by 1, horizontal sep=16mm}, omrecord, grid=none,
  xmin=0, xmax=35, ymin=0, ymax=35, width=0.46\linewidth, height=0.42\linewidth, samples=80,
  xlabel={Labour $L$}, ylabel={Capital $K$}]
\nextgroupplot[title={\small Returns to scale on a ray}, width=0.46\linewidth]
  \addplot[demand, domain=2.9:35] {100/x};
  \addplot[demand, domain=11.5:35] {400/x};
  \addplot[demand, domain=25.8:35] {900/x};
  \addplot[shifted, domain=0.7:35] {21.54/x};
  \addplot[shifted, domain=1.6:35] {54.29/x};
  \addplot[shifted, domain=2.7:35] {93.22/x};
  \addplot[guide, domain=0:35] {x};
  \addplot[only marks, mark=*, black] coordinates {(10,10) (20,20) (30,30)};
  \addplot[only marks, mark=o, black] coordinates {(4.64,4.64) (7.37,7.37) (9.65,9.65)};
\nextgroupplot[title={\small Technical progress}, width=0.46\linewidth, xmax=20, ymax=20]
  \addplot[demand, domain=3.6:20, restrict y to domain=0:40] {(2*10^(-2/3) - x^(-2/3))^(-1.5)};
  \addplot[shifted, domain=2.86:20, restrict y to domain=0:40] {(2*8^(-2/3) - x^(-2/3))^(-1.5)};
  \addplot[supply, domain=2.4:20, restrict y to domain=0:40] {(2*10^(-2/3) - (1.5*x)^(-2/3))^(-1.5)};
\end{groupplot}
\end{tikzpicture}

{\small Left: \omterm{def:b1:technology:isoquant}{isoquants} for outputs 10, 20 and 30 of $(KL)^{1/2}$ (solid,
constant returns: evenly spaced on the ray $K = L$) and of $(KL)^{3/4}$
(dashed, increasing returns: closer and closer). Right: the \omterm{def:b1:technology:isoquant}{isoquant}
$Y = 10$ of a CES technology with $\sigma = 0.6$ (blue), after Hicks-neutral
progress of $25\%$ (dashed: it moves in towards the origin along every ray)
and after labour-augmenting progress of $50\%$ (red: every point moves left,
to two-thirds of its labour, which shifts the steep part of the \omterm{def:b1:technology:isoquant}{isoquant} a
long way and its flat part hardly at all).}
\end{omfigure}

\section{The elasticity of substitution}\label{sec:b1:technology:sigma}

The \omterm{def:b1:technology:mrts}{MRTS} falls as labour replaces capital along a convex \omterm{def:b1:technology:isoquant}{isoquant}. How
fast it falls says how easy the replacement is: when it falls slowly, a
small change in relative prices moves the firm far along the \omterm{def:b1:technology:isoquant}{isoquant}.

\begin{definition}[Elasticity of substitution]\label{def:b1:technology:sigma}
\index{elasticity of substitution}
The \emph{elasticity of substitution} at a point of an \omterm{def:b1:technology:isoquant}{isoquant} is the
percentage change in the capital--labour ratio per percentage change in the
\omterm{def:b1:technology:mrts}{MRTS}, along the \omterm{def:b1:technology:isoquant}{isoquant}:
\[
\sigma = \frac{\dd\ln(K/L)}{\dd\ln\mathrm{MRTS}}.
\]
\end{definition}

For Cobb--Douglas, $\mathrm{MRTS} = (\beta/\alpha)(K/L)$, so $\ln(K/L)$ and
$\ln\mathrm{MRTS}$ differ by a constant and $\sigma = 1$. For fixed proportions
the ratio $K/L$ never changes at efficient points, $\sigma = 0$; for \omterm{def:b1:preferences-and-utility:perfect}{perfect
substitutes} $\sigma$ is infinite. One family spans the range.

\begin{definition}[CES production function]\label{def:b1:technology:ces}
\index{CES production function}
A \emph{constant-elasticity-of-substitution} (CES) \omterm{def:b1:technology:production-function}{production function} is
\[
f(K, L) = A\bigl(\alpha K^\rho + (1 - \alpha)L^\rho\bigr)^{\nu/\rho},\qquad \rho < 1,\ \rho \ne 0,
\]
with $A > 0$, $0 < \alpha < 1$ and $\nu > 0$, the degree of homogeneity.
\end{definition}

\begin{proposition}[The CES elasticity]\label{prop:b1:technology:ces-sigma}
The CES function has $\sigma = 1/(1 - \rho)$ at every point. It tends to the
Cobb--Douglas function as $\rho \to 0$, to fixed proportions as $\rho \to -\infty$
and to \omterm{def:b1:preferences-and-utility:perfect}{perfect substitutes} as $\rho \to 1$.
\end{proposition}

\begin{proof}
The common factor $A\nu(\cdots)^{\nu/\rho - 1}$ cancels in the ratio of \omterm{def:b1:technology:marginal-product}{marginal
products}:
$\mathrm{MRTS} = f_L/f_K = \bigl((1 - \alpha)/\alpha\bigr)(K/L)^{1-\rho}$. So
$\ln\mathrm{MRTS} = \ln\bigl((1 - \alpha)/\alpha\bigr) + (1 - \rho)\ln(K/L)$, and
$\sigma = \dd\ln(K/L)/\dd\ln\mathrm{MRTS} = 1/(1 - \rho)$. The limits of $\sigma$ are 1, 0 and
infinity; the limit of the function as $\rho \to 0$ is
\cref{exo:b1:technology:cd-limit}.
\end{proof}

A firm that minimises cost sets the \omterm{def:b1:technology:mrts}{MRTS} equal to the ratio of the wage $w$
to the rental of capital $r$ (\cref{ch:b1:cost-minimisation}). So $\sigma$ is
also the elasticity of the chosen $K/L$ with respect to $w/r$: with
$\sigma = 0.5$, a doubling of the wage--rental ratio raises capital per worker
by $2^{0.5} - 1 = 41\%$; with $\sigma = 1$ it doubles it. In the opening chart,
if all countries had one CES technology and paid one rental for capital,
the slope would estimate $\sigma$. They do not, and the slope of $0.94$ mixes
substitution with differences in technology and in the price of capital.

\begin{omchart}{The capital--labour ratio against the wage--rental ratio,
Korea and the United States, 1953--2019 (Korea from 1970, its first observed
labour share), both indexed to 1980 = 1,
logarithmic axes; the dotted line has slope 1.}{Penn World Table 11.0:
capital stock at constant national prices per hour worked; wage--rental
ratio $= \text{labour share}/(1 - \text{labour share}) \times K/L$}
\begin{tikzpicture}
\begin{loglogaxis}[omrecord, log ticks with fixed point, xtick={0.5,1,2,5,10}, ytick={0.5,1,2,5,10}, xmin=0.3, xmax=20, ymin=0.3, ymax=20,
  xlabel={Wage--rental ratio (1980 = 1)}, ylabel={Capital per hour (1980 = 1)},
  width=0.6\linewidth, height=0.5\linewidth, legend pos=north west]
  \addplot[demand, mark=none] table[x=wr_kor, y=kl_kor, col sep=comma]
    {figdata/bachelor-1/09-technology/out/kl_wr_paths.csv};
  \addplot[supply, mark=none] table[x=wr_usa, y=kl_usa, col sep=comma]
    {figdata/bachelor-1/09-technology/out/kl_wr_paths.csv};
  \addplot[guide, domain=0.3:20, forget plot] {x};
  \legend{Korea, United States}
\end{loglogaxis}
\end{tikzpicture}
\end{omchart}

The rental of capital is not observed directly. The chart backs it out:
with constant returns and inputs paid their \omterm{def:b1:technology:marginal-product}{marginal products}, the ratio
of the wage bill to the capital income, $wL/(rK)$, is the labour share over
its complement, so $w/r$ is that ratio times $K/L$. From 1980 to 2019 Korea's
capital per hour rose 12.9 times and its wage--rental ratio 15.6 times: a
slope of $0.93$. In the United States the two rose 1.66 and 1.45 times, a
slope of $1.37$. % figdata: kl_wr_slopes
Read as estimates of $\sigma$, these numbers assume that the technology stood
still; it did not, and the next section shows how \omterm{def:b1:technology:technical-progress}{technical progress} moves
both ratios.

\begin{record}[The elasticity of substitution]
Seventy-five years of estimates of the \omterm{def:b1:technology:sigma}{elasticity of substitution} between
capital and labour, from time series of industries, panels of firms and
cross-sections of countries, vary widely. A survey of them concludes that
the weight of the evidence suggests a value of $\sigma$ between $0.40$ and
$0.60$, and that there is little evidence for the Cobb--Douglas function's
$\sigma = 1$ (Chirinko 2008). % ledger: b1:chirinko2008
Short-run data understate the long-run elasticity, because firms take
years to change their machines; regressions over long periods mix
substitution with technical change. Both problems are visible in the
second chart.
\end{record}

\begin{model}{The Cobb--Douglas technology}
\assumes $Y = AK^\alpha L^{1-\alpha}$ with $0 < \alpha < 1$: constant returns and an \omterm{def:b1:technology:sigma}{elasticity
of substitution} of one.
\predicts Under cost minimisation the shares of capital and labour in costs
are $\alpha$ and $1 - \alpha$ whatever the wage and the rental (\cref{ch:b1:cost-minimisation});
across economies with one rental of capital, capital per hour rises in
proportion to the wage.
\tested Consistent with it, but not a test: across 112 countries in 2019,
capital per hour rises with the wage with a slope of $0.94$ in logarithms,
close to one, though the slope also mixes in differences in technology and
in the price of capital. % figdata: capital_wage_stats
Factor shares were roughly stable for decades in many countries
(\cref{ch:b1:factor-markets}).
\breaks Labour shares in the same data range from 0.33 to 0.63 (tenth to
ninetieth percentile), % figdata: capital_wage_stats
which one Cobb--Douglas technology with one $\alpha$ cannot produce. Most firm and industry estimates put $\sigma$ between $0.4$ and $0.6$
(Chirinko 2008); % ledger: b1:chirinko2008
labour shares have fallen in many countries since 1980
(\cref{ch:b1:factor-markets}), which one Cobb--Douglas function cannot
produce.
\end{model}

\begin{measure}[A capital stock]
A country's capital stock is not counted; it is built up from investment by
the \emph{perpetual inventory method}: this year's stock of an asset is
last year's, minus depreciation, plus this year's investment at constant
prices, $K_t = (1 - \delta)K_{t-1} + I_t$. The Penn World Table does this for up
to six kinds of asset, each with its own geometric depreciation rate,
assumed the same in every country and every year: $2\%$ a year for
structures, $11.5\%$ for communication equipment, $12.6\%$ for other
machinery, $18.9\%$ for transport equipment and $31.5\%$ for computers and
software (Inklaar and Timmer 2013). % ledger: b1:pwt-capital
A country that invests more in computers therefore has a capital stock
that wears out faster; and a stock measured in money adds up ovens,
roads and software by their prices.
\end{measure}

\section{Technical progress}\label{sec:b1:technology:progress}

A better recipe, a faster oven or a better-trained baker changes the
\omterm{def:b1:technology:production-function}{production function} itself.

\begin{definition}[Technical progress]\label{def:b1:technology:technical-progress}
\index{technical progress}\index{Hicks-neutral technical progress}\index{labour-augmenting technical progress}\index{capital-augmenting technical progress}
\emph{Technical progress} is a change in the \omterm{def:b1:technology:production-function}{production function} that
raises the output of some input bundles and lowers that of none. It is
\emph{Hicks-neutral} if $f$ becomes $Af$ with $A > 1$, leaving every \omterm{def:b1:technology:mrts}{MRTS}
unchanged; \emph{labour-augmenting} if $f(K, L)$ becomes $f(K, bL)$ with
$b > 1$, as if each hour did the work of $b$ hours; and
\emph{capital-augmenting} if it becomes $f(bK, L)$.
\end{definition}

For Cobb--Douglas the three are the same thing: $f(K, bL) = b^{1-\alpha}f(K, L)$.
For other technologies they differ, and the difference shows in factor
shares. With $\sigma < 1$, labour-augmenting progress is like a fall in the
price of labour per effective hour; since inputs substitute badly, the firm
raises its effective hours per unit of capital by less than their price
falls (and uses fewer physical hours per unit of capital), so the cost
share of capital \emph{rises} (\cref{exo:b1:technology:augmenting}). A measured
productivity rise therefore says that output per unit of input went up,
not which input became more productive; and a change in factor shares can
come from technology as well as from prices.

\section*{Exercises}\label{sec:b1:technology:exercises}

\begin{exercise}[$\star$]\label{exo:b1:technology:mrts-cd}
Compute the \omterm{def:b1:technology:mrts}{MRTS} of $f = K^{0.3}L^{0.7}$ at $(K, L) = (8, 27)$, and say what it means.
\end{exercise}

\begin{exercise}[$\star$]\label{exo:b1:technology:three}
Give the \omterm{def:b1:technology:returns-to-scale}{returns to scale} of $f_1 = K^{0.5}L^{0.5}$, $f_2 = K^{0.6}L^{0.6}$ and
$f_3 = \bigl(\min(K, 2L)\bigr)^{0.8}$.
\end{exercise}

\begin{exercise}[$\star$]\label{exo:b1:technology:read}
From the opening chart's table, an hour of work was paid 45.86 dollars and
used 256.58 dollars of capital in the United States, against 3.97 and
27.98 in India. % figdata: capital_wage_2019
Compute the two ratios. What would they be if both countries used one
Cobb--Douglas technology and paid one rental for capital?
\end{exercise}

\begin{exercise}[$\star$]\label{exo:b1:technology:shapes}
Each table gives three input bundles $(K, L)$ that make the same output.
Which shape does each technology look like: (a) $(1, 9)$, $(3, 3)$, $(9, 1)$;
(b) $(2, 8)$, $(5, 5)$, $(8, 2)$; (c) $(2, 4)$, $(2, 9)$, $(7, 4)$?
\end{exercise}

\begin{exercise}[$\star\star$]\label{exo:b1:technology:ces-sigma}
Compute the \omterm{def:b1:technology:mrts}{MRTS} of $f = (0.3K^{-1} + 0.7L^{-1})^{-1}$ and its \omterm{def:b1:technology:sigma}{elasticity of
substitution}.
\end{exercise}

\begin{exercise}[$\star\star$]\label{exo:b1:technology:cd-limit}
Show that $\bigl(\alpha K^\rho + (1 - \alpha)L^\rho\bigr)^{1/\rho} \to K^\alpha L^{1-\alpha}$ as $\rho \to 0$. (Take
logarithms and use l'H\^opital's rule.)
\end{exercise}

\begin{exercise}[$\star\star$]\label{exo:b1:technology:predict}
The wage--rental ratio doubles. Predict the percentage change in a
cost-minimising firm's capital per worker for $\sigma = 0.5$ and $\sigma = 1$.
\end{exercise}

\begin{exercise}[$\star\star$]\label{exo:b1:technology:local}
Compute the elasticity of scale of $f = K^{0.5}L^{0.5} - 2$ at $(4, 4)$ and at
$(16, 16)$. Where are returns increasing?
\end{exercise}

\begin{exercise}[$\star\star$]\label{exo:b1:technology:estimate}
From the second chart's table, between 1980 and 2019 capital per hour rose
12.872 times in Korea and 1.656 times in the United States, and the
wage--rental ratio 15.61 and 1.445 times. % figdata: kl_wr_slopes
Estimate $\sigma$ for each. Why might these be poor estimates?
\end{exercise}

\begin{exercise}[$\star\star\star$]\label{exo:b1:technology:euler-proof}
Prove the converse of Euler's theorem: if $Kf_K + Lf_L = kf$ for all
$(K, L)$, then $f$ is homogeneous of degree $k$. (Study
$g(t) = t^{-k}f(tK, tL)$.)
\end{exercise}

\begin{exercise}[$\star\star\star$]\label{exo:b1:technology:augmenting}
A firm with the CES technology $\bigl(\alpha K^\rho + (1 - \alpha)(bL)^\rho\bigr)^{1/\rho}$ minimises
cost at given $w$ and $r$. Show that the ratio of its capital cost to its
labour cost is $\bigl(\alpha/(1 - \alpha)\bigr)^\sigma\bigl(w/(br)\bigr)^{\sigma - 1}$. Does the
capital share rise or fall when $b$ rises, for $\sigma < 1$?
\end{exercise}

\begin{exercise}[$\star\star\star$]\label{exo:b1:technology:refute}
What evidence in the opening chart's data would refute ``all these
countries use the same Cobb--Douglas technology and pay the same rental for
capital''? Is the slope of $0.94$ such evidence?
\end{exercise}

\section{Problem: Two Bakeries and a New Oven}\label{sec:b1:technology:problem}

\begin{problem}[{Weekend problem --- the elasticity of substitution revealed by two bakeries}]\label{pb:b1:technology:bakeries}
Two bakeries (invented) have the same technology,
\[
Y = 100\bigl(\tfrac12K^\rho + \tfrac12L^\rho\bigr)^{1/\rho},\qquad \rho = -\tfrac23,
\]
where $K$ is ovens, $L$ bakers and $Y$ loaves a day. Bakery A, in Kestria,
uses 2 ovens and 8 bakers; bakery B, in Brenland, 4 ovens and 4 bakers.

\medskip\noindent\textbf{Part I --- Isoquants.}
\begin{enumerate}
  \item Compute each bakery's output.
  \item Show that the \omterm{def:b1:technology:mrts}{MRTS} is $(K/L)^{5/3}$ and compute it for each bakery.
  \item How many ovens would bakery A need to bake B's 400 loaves with its 8 bakers?
  \item Show that no number of bakers can bake 400 loaves with fewer than $4 \times 2^{-3/2} = 1.414$ ovens. How does this \omterm{def:b1:technology:isoquant}{isoquant} differ from a Cobb--Douglas one?
  \item Compute bakery A's average products of bakers and of ovens.
  \item Compute bakery A's \omterm{def:b1:technology:marginal-product}{marginal products}, $f_L = \tfrac12 \cdot 100^\rho(Y/L)^{1-\rho}$ and the same for $K$, and check their ratio against question~2.
\end{enumerate}

\medskip\noindent\textbf{Part II --- Scale.}
\begin{enumerate}[resume]
  \item Compute the output of a bakery with twice A's inputs. What are the \omterm{def:b1:technology:returns-to-scale}{returns to scale}?
  \item Check Euler's theorem at bakery A.
  \item A larger bakery uses $Y = 100\bigl(\tfrac12K^\rho + \tfrac12L^\rho\bigr)^{1.2/\rho}$. Compute its output with B's inputs.
  \item Compute its output with half of B's inputs, and compare two such half-size bakeries with one at B's size.
  \item What does the exponent 1.2 mean?
  \item Why do bakeries stop growing, if the second technology holds?
\end{enumerate}

\medskip\noindent\textbf{Part III --- Substitution.} In Kestria an oven rents
for ten times a baker's daily wage, $w/r = 0.1$; in Brenland they cost the
same, $w/r = 1$. Each bakery minimises cost.
\begin{enumerate}[resume]
  \item Compute $\sigma$ from $\rho$.
  \item Show that a cost-minimising bakery chooses $K/L = (w/r)^\sigma$.
  \item An economist sees only the two choices, $K/L = 1/4$ and $1$, and the two wage--rental ratios. What \omterm{def:b1:technology:sigma}{elasticity of substitution} do they reveal?
  \item Check bakery A's choice against the technology's prediction.
  \item Brenland's wage--rental ratio rises by $10\%$. Predict the change in bakery B's ovens per baker.
  \item Compute the share of ovens in each bakery's costs, $rK/(rK + wL)$, and explain why it is larger in Kestria.
\end{enumerate}

\medskip\noindent\textbf{Part IV --- A new oven.} A new oven bakes like 1.5 old
ones and rents for the same price: capital is augmented, $b = 1.5$.
\begin{enumerate}[resume]
  \item Compute bakery A's output with new ovens and its old inputs.
  \item What is the rental of one unit of \emph{effective} oven, $bK$? Compute the cost-minimising effective ovens per baker in each country.
  \item Compute the new physical ovens per baker in each country. Do the bakeries rent more or fewer ovens per baker?
  \item Compute the new oven cost shares.
  \item Explain why an improvement in ovens lowers the share of ovens in costs.
  \item State the result: the \omterm{def:b1:technology:sigma}{elasticity of substitution} revealed by the two bakeries.
\end{enumerate}
\end{problem}
